\documentclass[12pt]{article}
\usepackage{mgates-letter}
\definecolor{dark_blue} {rgb}{0., 0., 0.65}

\usepackage{textcomp}
\usepackage{mathrsfs}  % mathscr font
\usepackage{boxedminipage}
\usepackage{rotating}
\usepackage[inline]{enumitem}
%\usepackage{natbib}
\usepackage{xcolor}
\usepackage[colorlinks, filecolor=dark_blue, urlcolor=dark_blue, linkcolor=black, citecolor=black]{hyperref}

\newcommand{\todoinline}[1]{{\color{violet} #1}}

\begin{document}

\begin{titlepage}

	\newcommand{\HRule}{\rule{\linewidth}{0.5mm}}
	\center

	\textsc{\Large Ph.D. Programme in Computer Science And Engineering}\\[0.5cm]

	\textsc{\Large XL Cycle}\\[0.6cm]

	\hrule width \hsize \kern 1mm \hrule width \hsize height 2pt
	\vspace{0.8cm}
	{ \large \bfseries Ph.D. Period Abroad Final Report for Marco Polo}\\[0.6cm]
	{ \large Integration of Control Barrier Functions for Safety-Critical Guarantees in Aggregate Computing}\\[0.6cm]

	\bfseries{May, 2026}

    \vspace{1.5cm}

    \noindent
    \begin{minipage}[t]{0.3\textwidth}
        \raggedright
        \textbf{Supervisors:}\\[0.5cm]
        Prof. Danilo Pianini\\
        Prof. Mirko Viroli\\
        Prof. Enrico Gallinucci
    \end{minipage}%
	\hfill
    \begin{minipage}[t]{0.3\textwidth}
        \centering
        \textbf{Abroad Supervisor:}\\[0.5cm]
        Alessandro Vittorio Papadopoulos
    \end{minipage}
    \hfill
    \begin{minipage}[t]{0.3\textwidth}
        \raggedleft
        \textbf{PhD Candidate:}\\[0.5cm]
        Angela Cortecchia
    \end{minipage} \\[0.6cm]

	\hrule width \hsize height 2pt \kern 1mm \hrule width \hsize height 1pt
	\vspace{0.4cm}

\end{titlepage}

\section{The stay}
During my research stay abroad,
carried out at Mälardalen University (MDU) from January to April 2026,
I collaborated with the MARC-Mälardalen University Automation Research Center,
under the supervision of Prof. Alessandro Vittorio Papadopoulos.
%
The main objective of the stay was to investigate how to integrate Control Lyapunov Functions and Control Barrier Functions into Aggregate Computing,
so that collective adaptive systems can guarantee stability and safety during transient adaptation phases,
not only eventual consistency.

This period at MDU represented a valuable opportunity to align my PhD research with the expertise and
 ongoing work of the host group.
%
The collaboration was fully consistent with the objectives outlined in my period abroad proposal and
 enabled the exploration of new methodological perspectives,
 as well as the establishment of promising research synergies.

\section{Research context and objectives}
The research activity conducted during the stay was situated within the broader context
of Collective Adaptive Systems and Aggregate Computing,
focusing on the integration of formal control-theoretic guarantees into collective
adaptive behaviors.

In particular,
the goal was to investigate how Control Lyapunov Functions and Control Barrier Functions
can be integrated into the Aggregate Computing framework,
enabling collective systems to reason not only about convergence and stability,
but also about the enforcement of safety constraints during transient adaptation phases.
%
This is especially relevant for safety-critical scenarios,
such as multi-robot systems,
where eventual consistency alone is not sufficient to prevent unsafe behaviors while
the system is reorganizing toward a desired collective configuration.
%
The interaction with the host research group also allowed these topics to be analyzed
from a control-theoretic perspective,
highlighting the challenges involved in translating collective-level goals into
local safety-preserving actions.

\section{Main activities}
During my stay at Mälardalen University, I carried out the following main activities:
\begin{itemize}
    \item weekly research meetings with Prof. Alessandro Vittorio Papadopoulos,
    aimed at discussing the integration of control-theoretic methods into Aggregate Computing;
    \item in-depth study and analysis of the relevant literature on Control Lyapunov Functions,
    Control Barrier Functions, and their use for safety-critical guarantees in distributed and
    adaptive systems;
    \item investigation of possible ways to map collective-level objectives and safety requirements
    into local control constraints enforceable by individual devices;
    \item preliminary definition of safety-critical use cases, with particular attention to
    multi-robot systems, to assess how the proposed integration could guarantee safety during
    transient adaptation phases;
    \item regular meetings with the MARC research group to present ongoing work.
\end{itemize}

\section{Ongoing and future work}
By the end of the stay,
the main objectives of the research activity had been achieved.
%
In particular,
we obtained a working integration of Control Lyapunov Functions and Control Barrier Functions
into the Aggregate Computing framework,
and validated the approach through simulation.

The results obtained during the stay are currently being consolidated into a first scientific publication,
which we aim to complete in the coming weeks.
%
At the same time,
the collaboration with the host research group is continuing,
with the objective of extending the work and preparing further publications.
%
Future developments will focus in particular on more advanced simulation scenarios,
including the definition and control of actual drone formations,
in order to assess the proposed approach in increasingly realistic safety-critical settings.


%\bibliographystyle{unsrt}
%\bibliography{bibliography}

\end{document}
